\subsection{Applications of functional calculus}

PROPOSITION 9. --- Every injective morphism of stellar algebras is isometric and, in particular, has closed image.

Let A and A′ be stellar algebras, and let \(\pi\colon A\to A'\) be a morphism of involutive algebras from A to A'.

Suppose first that A and A' are unital and that \(\pi\) is unital. We have \(\|\pi\| \leqslant 1\) (prop. 2). Suppose there exists x in A such that \(\|\pi(x)\| < \|x\|\) . Let \(y = x^{*}x\) ; it is a Hermitian element of A. Since A and A' are stellar algebras, we have \(\|\pi(y)\| = \|\pi(x)\|^{2} < \|x\|^{2} = \|y\|\) , that is, \(\varrho(\pi(y)) < \varrho(y)\) (lemma 7 of I, p.~104). In particular, \(\mathrm{Sp}_{\mathrm{A}'}(\pi(y))\) is a closed subset of \(\mathrm{Sp}_{\mathrm{A}}(y)\) , distinct from \(\mathrm{Sp}_{\mathrm{A}}(y)\)

(Remark 6 of I, p.~3 and Th. 1 of I, p.~24). There then exists a nonzero function \(f \in \mathcal{C}(\mathrm{Sp}_{\mathrm{A}}(y))\) such that \(f| \mathrm{Sp}_{\mathrm{A}'}(\pi(y)) = 0\) (TG, IX, p.~13, Prop. 3). Let \(w = f(y) \in \mathrm{A}\) . We have \(w \neq 0\) since \(f \neq 0\) but \(\pi(w) = \pi(f(y)) = f(\pi(y)) = 0\) since \(f\) is zero on \(\mathrm{Sp}_{\mathrm{A}'}(\pi(y))\) (Prop. 8). Hence \(\pi\) is not injective.

We now treat the general case. Let \(\widetilde{\mathrm{A}}\) and \(\widetilde{\mathrm{A}}'\) be the stellar algebras deduced from A and from \(\mathrm{A}'\) respectively by adjoining a unit element (def. 4 of I, p.~106). There exists a unique unital morphism of involutive algebras \(\widetilde{\pi}\colon \widetilde{\mathrm{A}}\to \widetilde{\mathrm{A}}'\) extending \(\pi\) . This morphism is injective, hence is isometric by what precedes. For every \(x\in \mathrm{A}\) , we then have \(\| \pi (x)\| = \| \widetilde{\pi} (x)\| = \| x\|\) .

Lemma 11. --- Let X be a completely regular topological space, that is, uniformizable and separated (TG, IX, p.~8, def. 4), containing at least two points. There exist nonzero continuous functions f and g in \(\mathcal{C}(\mathrm{X})\) such that \(fg = 0\) .

Let \(x \neq y\) be distinct points of X. Let U and V be open neighborhoods of x and y, respectively, such that U and V are disjoint. Since X is uniformizable, by TG, IX, p.~7, th. 2, there exists a function \(f \in \mathcal{C}(X)\) such that \(f(x) = 1\) and \(f|X - U = 0\) . Similarly, there exists \(g \in \mathcal{C}(X)\) such that \(g(y) = 1\) and \(g|X - V = 0\) . We then have fg = 0.

PROPOSITION 10. --- Let A be a unital stellar algebra. Suppose that for every pair \((x,y)\) of permutable elements of A, the condition xy = 0 implies that x = 0 or y = 0. Then A = C · 1.

If A is not equal to C·1, there exists a hermitian element x in A which does not belong to C·1 (lemma 2 of I, p.~96). Let B be the unital stellar subalgebra of A generated by x. It is commutative, and isomorphic to \(\mathcal{C}(\mathrm{Sp}_{\mathrm{A}}(x))\) (I, p.~110, remark). Since x is not scalar, its spectrum in B is not reduced to a single element (example 1 of I, p.~110). There therefore exist nonzero continuous functions f and g on \(\mathrm{Sp}_{\mathrm{A}}(x)\) such that fg = 0 (lemma 11). The elements \(f(x)\) and \(g(x)\) of A are nonzero, permutable, and satisfy \(f(x)g(x) = 0\) in A.

PROPOSITION 11. --- Let A be a unital stellar algebra, and let a, x and y be elements of A. Suppose that x and y are normal. If xa = ay, then we have \(f(x)a = af(y)\) for every continuous function f on the union of the spectrum of x and the spectrum of y. In particular, we have \(f(aa^{*})a = af(a^{*}a)\) for every function \(f \in \mathcal{C}(\mathrm{Sp}'(a^{*}a))\) .

Let \(\mathrm{S} = \mathrm{Sp}(x) \cup \mathrm{Sp}(y)\) . Proposition 6 of I, p.~106 implies that \(x^{*}a = ay^{*}\) . Consequently, we obtain \(f(x)a = af(y)\) for every function \(f\) which is of the form \(z \mapsto \mathrm{P}(z,\overline{z})\) , where \(\mathrm{P} \in \mathbf{C}[\mathrm{X},\mathrm{Y}]\) is a polynomial. Since the set of functions \(f \in \mathcal{C}(\mathrm{S})\) satisfying \(f(x)a = af(y)\) is a closed subalgebra of \(\mathcal{C}(\mathrm{S})\) , it coincides with \(\mathcal{C}(\mathrm{S})\) by TG, X, p.~40, cor. 1.

The second assertion is a consequence of the first, applied to the hermitian elements \(x = aa^{*}\) and \(y = a^{*}a\) , taking into account the fact that \(\mathrm{Sp}'(a^{*}a) = \mathrm{Sp}'(aa^{*})\) (prop. 1 of I, p.~5).

